
This project set out to understand when and why goodness-of-fit testing becomes hard in
high-dimensional models, and how modern machine-learning-inspired procedures compare to
classical statistical tests. The Gaussian sequence model over Sobolev ellipsoids provided
a clean and mathematically tractable setting in which these questions could be made precise.

Chapter~3 revisited the classical theory of minimax distinguishability through the lens of
$\chi^2$-type tests. The key takeaway is that the smoothness of the signal class completely
governs the difficulty of the problem: there is a sharp phase transition below which no
test can do better than random guessing, and above which consistent testing becomes
possible. This dichotomy is a fundamental feature of the problem, independent of the
specific test used.

Chapter~4 then asked whether a classifier-based approach could achieve the same limits.
The answer, perhaps surprisingly, is yes. We proved that any test requires at least
\[
    n(\varepsilon, \delta, s)
    \;\gtrsim\;
    \frac{\sqrt{\log(1/\delta)}}{\varepsilon^{\frac{4s+1}{2s}}}
    + \frac{\log(1/\delta)}{\varepsilon^{2}}
\]
observations, and that the CAT test based on a truncated likelihood-ratio separating set
matches this bound. The two terms in the lower bound tell different stories: the first
captures the nonparametric hardness of the problem through the smoothness parameter $s$,
while the second is simply the unavoidable cost of achieving high confidence $\delta$ in
any binary testing problem.

Of course, CAT is not the only way to approach goodness-of-fit testing. Classical tests
such as Kolmogorov--Smirnov, $\chi^2$ goodness-of-fit, or kernel-based methods like the
Maximum Mean Discrepancy have been studied for decades and remain widely used in practice.
What makes CAT interesting is not that it outperforms these methods, but that it achieves
the same fundamental limits while being naturally suited to the modern setting where
distributions are only accessible through simulators.

More broadly, the results presented here suggest that the boundary between statistical
theory and machine learning is less sharp than it might seem: a procedure motivated by
classification can, with the right analysis, be shown to be optimal in a rigorous
minimax sense.